\begin{definition}[Regrouping original physical legs by torus bonds]%
    \label{def:peps_torus_physical_bond_regrouping}
    \lean{TNLean.PEPS.torusSiteBondEndpointEquiv,
        TNLean.PEPS.torusSiteBondEndpointEquiv_apply_horizontal,
        TNLean.PEPS.torusSiteBondEndpointEquiv_apply_vertical,
        TNLean.PEPS.torusSiteBondEndpointEquiv_symm_apply,
        TNLean.PEPS.torusBondRegrouping,
        TNLean.PEPS.torusBondRegrouping_apply,
        TNLean.PEPS.torusBondRegroupingIsometry,
        TNLean.PEPS.torusBondRegroupingIsometry_apply}
    \leanok
    \uses{lem:peps_torus_projector_expansion}
    Let $X$ be the vertex set of a torus with positive periods and let
    $\mathcal B=X\times\{\mathrm h,\mathrm v\}$ be its horizontal and
    vertical bond positions. Write $\sigma_v=(t_v,r_v,b_v,l_v)$ for the
    four original physical coordinates at $v$. Their bond regrouping is
    \begin{align}
        \beta_{v,\mathrm h} & =(l_{v+\mathrm{east}},r_v),  &
        \beta_{v,\mathrm v} & =(t_v,b_{v+\mathrm{north}}).
        \label{eq:peps_physical_bond_regrouping}
    \end{align}
    Each pair is ordered head, tail: horizontal bonds point right and
    vertical bonds point down. The inverse reads
    \begin{align}
        t_v & =\beta^1_{v,\mathrm v},                & r_v & =\beta^2_{v,\mathrm h}, &
        b_v & =\beta^2_{v-\mathrm{north},\mathrm v}, &
        l_v & =\beta^1_{v-\mathrm{east},\mathrm h}.
    \end{align}
    Thus regrouping is a bijection of the actual physical basis
    configurations. Its induced coefficient map is
    $(\mathcal R\psi)(\beta)=\psi(\sigma(\beta))$, and is a linear
    isometric equivalence of the finite-dimensional physical Hilbert spaces.
    This is the grouping into the two physical bond endpoints used in
    \cite[Section~7]{Schuch2010PEPS}.
\end{definition}
\begin{proof}\leanok
    First use the torus bond-to-site coordinate bijection, then pair the
    head and tail coordinates of each horizontal and vertical bond.
    The displayed inverse recovers all four site coordinates. A bijection
    of orthonormal basis configurations induces a linear isometry.
\end{proof}

\begin{lemma}[Overlap and product identities under bond regrouping]%
    \label{lem:peps_torus_physical_bond_regrouping_overlap}
    \lean{TNLean.PEPS.torusBondRegrouping_dotProduct,
        TNLean.PEPS.torusBond_product}
    \leanok
    \uses{def:peps_torus_physical_bond_regrouping}
    Every pair of physical coefficient vectors satisfies
    \begin{align}
        \langle\mathcal R\psi,\mathcal R\phi\rangle
         & =\langle\psi,\phi\rangle.
    \end{align}
    For any complex bond factors $F_e$, the literal product identity is
    \begin{align}
        \prod_{v\in X}F_{v,\mathrm h}F_{v,\mathrm v}
         & =\prod_{e\in\mathcal B}F_e.
    \end{align}
\end{lemma}
\begin{proof}\leanok
    Reindex the finite overlap sum by the physical configuration bijection.
    For the product, separate the vertex and direction indices and use
    commutativity of the scalar factors.
\end{proof}

\begin{theorem}[Actual torus bond coefficients after physical regrouping]%
    \label{thm:peps_torus_regrouped_averaging_bonds}
    \lean{TNLean.PEPS.torusBondRegrouping_averagingSite,
        TNLean.PEPS.torusBondRegrouping_dressedAveragingSite,
        TNLean.PEPS.torusBondRegrouping_thetaWeightedAveragingSite}
    \leanok
    \uses{def:peps_torus_physical_bond_regrouping,
        lem:peps_torus_physical_bond_regrouping_overlap,
        lem:peps_torus_projector_expansion,
        thm:peps_theta_weighted_torus_bonds}
    Let $U$ be a finite-group representation on the virtual space, and
    let $W$ commute with every $U(g)$. Denote by $\Psi_{U,W}$ the actual
    torus state of the normalized averaging site with $W$ on each virtual
    leg. Put $N=|X|$, and let $h(e),t(e)$ be the head and tail of a bond
    in the orientations of (\ref{eq:peps_physical_bond_regrouping}). Then
    \begin{align}
        (\mathcal R\Psi_{U,W})(\beta)
         & =|G|^{-N}\sum_{q:X\to G}\prod_{e\in\mathcal B}
        \bigl[W^2U(q_{h(e)}q_{t(e)}^{-1})\bigr]_
        {\beta_e^1,\beta_e^2}.
        \label{eq:peps_actual_regrouped_weighted_bonds}
    \end{align}
    In particular, $W=\Theta$ gives the actual square-root-weighted
    bond vectors in \cite[Section~7]{Schuch2010PEPS}.
    For ordinary averaging sites with inserted bond matrices $O_e$, the
    regrouped coefficients instead contain
    $[U(q_{h(e)})O_eU(q_{t(e)}^{-1})]_{\beta_e^1,\beta_e^2}$.
    All coefficients are obtained from the original torus contraction;
    no equality of states is assumed.
\end{theorem}
\begin{proof}\leanok
    Apply the averaging-site expansion to the recovered site
    configuration. The inverse coordinate formulas identify each matrix
    row and column with its actual physical bond endpoint. For dressed
    sites, the two endpoint weights combine to $W^2$, which commutes
    through the relative group action. The product identity then groups
    the horizontal and vertical factors into the product over all bonds.
\end{proof}
